\section{Week 1: What Is a Cognitive Model?}

\begin{quote}
\textbf{Chapter 1: Fundamentals of cognitive modeling} · Lecture 1 h ·
Lab 2 h \textbf{This week's running model:} a simple random walk model
of two-choice decisions.
\end{quote}

\begin{center}\rule{0.5\linewidth}{0.5pt}\end{center}

\subsection{Key terms}

Important terms of this lesson, with their Persian equivalents (repeated
even if they appeared in earlier weeks).

\begin{longtable}[]{@{}
  >{\raggedright\arraybackslash}p{(\linewidth - 4\tabcolsep) * \real{0.3333}}
  >{\raggedright\arraybackslash}p{(\linewidth - 4\tabcolsep) * \real{0.3333}}
  >{\raggedright\arraybackslash}p{(\linewidth - 4\tabcolsep) * \real{0.3333}}@{}}
\caption{\centering Key terms}\tabularnewline
\toprule\noalign{}
\begin{minipage}[b]{\linewidth}\raggedright
English
\end{minipage} & \begin{minipage}[b]{\linewidth}\raggedright
\RL{\CMfa فارسی}
\end{minipage} & \begin{minipage}[b]{\linewidth}\raggedright
\RL{\CMfa توضیح کوتاه}
\end{minipage} \\
\midrule\noalign{}
\endfirsthead
\toprule\noalign{}
\begin{minipage}[b]{\linewidth}\raggedright
English
\end{minipage} & \begin{minipage}[b]{\linewidth}\raggedright
\RL{\CMfa فارسی}
\end{minipage} & \begin{minipage}[b]{\linewidth}\raggedright
\RL{\CMfa توضیح کوتاه}
\end{minipage} \\
\midrule\noalign{}
\endhead
\bottomrule\noalign{}
\endlastfoot
cognitive model & \RL{\CMfa مدل شناختی} & \RL{\CMfa توصیف دقیق و ریاضی یا محاسباتی
از یک فرایند ذهنی} \\
random walk & \RL{\CMfa گام تصادفی} & \RL{\CMfa جمع شدن گام‌به‌گام شواهد تا رسیدن به
یک مرز} \\
verbal theory & \RL{\CMfa نظریه‌ی کلامی} & \RL{\CMfa نظریه‌ای که فقط با کلمات بیان شده
و پیش‌بینی دقیق نمی‌دهد} \\
Marr's levels & \RL{\CMfa سطوح تحلیل} Marr &
\RL{\CMfa سطح محاسباتی، الگوریتمی و تحلیل} \\
parameter & \RL{\CMfa پارامتر} & \RL{\CMfa کمیت قابل تنظیم مدل که بین افراد یا شرایط
فرق می‌کند} \\
noise & \RL{\CMfa نویز} & \RL{\CMfa تغییرات تصادفی در شواهد یا در داده} \\
simulation & \RL{\CMfa شبیه‌سازی} & \RL{\CMfa تولید داده‌ی مصنوعی با اجرای مدل} \\
drift rate & \RL{\CMfa نرخ رانش} & \RL{\CMfa سرعت میانگین جمع شدن شواهد؛ نشانه‌ی
کیفیت شواهد} \\
boundary separation & \RL{\CMfa فاصله‌ی مرزها} & \RL{\CMfa مقدار شواهد لازم برای
پاسخ؛ نشانه‌ی احتیاط} \\
evidence & \RL{\CMfa شواهد} & \RL{\CMfa اطلاعاتی که برای تصمیم جمع می‌شود} \\
accuracy & \RL{\CMfa دقت} & \RL{\CMfa نسبت پاسخ‌های درست} \\
response time (RT) & \RL{\CMfa زمان واکنش} & \RL{\CMfa فاصله‌ی ارائه‌ی محرک تا
پاسخ} \\
formal model / formalize & \RL{\CMfa مدل فرمال} / \RL{\CMfa فرمال کردن} & \RL{\CMfa نوشتن
نظریه به شکل معادله یا الگوریتم} \\
fit / fitting & \RL{\CMfa برازش} & \RL{\CMfa تنظیم پارامترها تا مدل با داده جور
شود} \\
forward / inverse problem & \RL{\CMfa مسئله‌ی مستقیم} / \RL{\CMfa معکوس} &
\RL{\CMfa شبیه‌سازی (پارامتر} \ensuremath{\leftarrow} \RL{\CMfa داده) در برابر برآورد (داده} \ensuremath{\leftarrow}
\RL{\CMfa پارامتر)} \\
non-decision time & \RL{\CMfa زمان غیرتصمیمی} & \RL{\CMfa زمان دیدن محرک و اجرای
حرکت پاسخ} \\
speed--accuracy trade-off & \RL{\CMfa بده‌بستان سرعت و دقت} & \RL{\CMfa سریع‌تر پاسخ
دادن به قیمت خطای بیشتر} \\
trial & \RL{\CMfa کوشش} (trial) & \RL{\CMfa یک بار انجام تکلیف توسط شرکت‌کننده} \\
temperature & \RL{\CMfa دما} (temperature) & \RL{\CMfa پارامتر نویز در انتخاب}
token\RL{\CMfa ؛} \ensuremath{\beta} = 1/T \\
random seed & seed \RL{\CMfa تصادفی} & \RL{\CMfa عدد شروع مولد اعداد تصادفی برای
تکرارپذیری} \\
quantile & \RL{\CMfa چندک} (quantile) & \RL{\CMfa مقداری که درصد مشخصی از داده زیر
آن است} \\
large language model (LLM) & \RL{\CMfa مدل بزرگ زبانی} & \RL{\CMfa شبکه‌ی عصبی‌ای که
کلمه‌ی بعدی متن را پیش‌بینی می‌کند} \\
token & token \RL{\CMfa (واحد متن)} & \RL{\CMfa کلمه یا بخشی از کلمه که مدل زبانی
پیش‌بینی می‌کند} \\
\end{longtable}

\begin{center}\rule{0.5\linewidth}{0.5pt}\end{center}

\subsection{Learning outcomes}

By the end of this week, students will be able to:

\begin{enumerate}
\def\labelenumi{\arabic{enumi}.}
\tightlist
\item
  \textbf{Define} a cognitive model and \textbf{distinguish} it from a
  verbal theory and from a purely statistical model of the data.
\item
  \textbf{Classify} models by purpose (descriptive, process, normative)
  and by family (symbolic, connectionist, stochastic/dynamical,
  Bayesian), and \textbf{place} them on Marr's three levels.
\item
  \textbf{Explain} how free parameters, a deterministic core and a noise
  model together turn a model into a \emph{data-generating} machine.
\item
  \textbf{Interpret} how changes in the model's parameters (drift,
  boundary) map onto psychological constructs (evidence quality,
  caution).
\end{enumerate}

\begin{center}\rule{0.5\linewidth}{0.5pt}\end{center}

\subsection{Motivation \& big picture}

Consider a familiar verbal theory of simple decisions:

\begin{quote}
\emph{``When people decide, they gather evidence until they have enough,
then they respond.''}
\end{quote}

It sounds complete, but try to get a prediction out of it:

\begin{itemize}
\tightlist
\item
  If we tell participants to \emph{be accurate} rather than \emph{be
  fast}, how much slower should they get? Twice as slow? 50 ms slower?
\item
  Are \textbf{errors} faster or slower than correct responses, or
  equally fast?
\item
  Should the response time (RT) distribution be symmetric or skewed?
\item
  If a stimulus becomes harder to see, should RTs change in the same way
  as under an accuracy instruction?
\end{itemize}

The verbal theory \textbf{cannot answer any of these}. Each answer
depends on decisions the theory never made: What counts as ``evidence''?
Does it arrive continuously or in chunks? What does ``enough'' mean? Is
the threshold the same for both responses?

As soon as we write the theory down as a model, for example a
\textbf{random walk} that steps up or down until it reaches a boundary,
every one of those questions gets a definite answer. Some answers will
turn out to be \emph{wrong}, and that is useful: a model that can be
wrong is a model we can learn from. This is the central message of the
course:

\begin{quote}
Formal models force us to make our assumptions explicit, produce precise
predictions, and turn ``Is this theory right?'' into ``How well, and
why, does this model account for the data?''
\end{quote}

\paragraph{Where is this used?}

Models like the random walk of this week are not only classroom
examples. They are used every day in research and practice:

\begin{itemize}
\tightlist
\item
  \textbf{Clinical psychology and psychiatry.} Patients with depression,
  ADHD or anxiety are often ``slower'' in simple tasks. Decision models
  show \emph{why}: lower evidence quality, more caution, or slower motor
  responses. These three explanations call for different interventions.
\item
  \textbf{Ageing research.} Older adults respond more slowly. Fitting
  decision models showed that much of this slowing comes from more
  caution and slower motor execution, not from worse information
  processing.
\item
  \textbf{Neuroscience.} Neurons in some brain areas increase their
  firing rate until a threshold before a decision, much like a random
  walk reaching a boundary. Model parameters are used to link behavior
  and brain activity.
\item
  \textbf{Applied settings.} Eyewitness decisions, medical image reading
  and quality control all involve deciding under noisy evidence. The
  same model separates ``how good is the evidence'' from ``how careful
  is the decision maker''.
\item
  \textbf{Artificial intelligence.} Ideas from cognitive models
  (evidence accumulation, learning from errors, probabilistic inference)
  are used in AI systems, and AI models are now tested with the methods
  of this course (see the AI \& LLMs section).
\end{itemize}

\begin{center}\rule{0.5\linewidth}{0.5pt}\end{center}

\Needspace{11\baselineskip}

\subsection{Lecture notes}

\paragraph{In plain words}

A \emph{theory} says something in words (``people collect evidence until
they have enough''). A \emph{model} says the same thing so precisely
that a computer can run it. Once it can run, the model makes exact
predictions: how often people will be correct, how long they will take,
and how both change when we change the task. This week's model is the
simplest possible decision model. Imagine a token on a number line. Each
moment it moves one step up (a bit of evidence for the correct answer)
or one step down (a bit of evidence for the wrong answer). When it
reaches the top line, the person answers correctly; when it reaches the
bottom line, they make an error. Only three numbers control everything:
how strongly the steps lean upward (\textbf{drift}, the quality of the
evidence), how far apart the two lines are (\textbf{boundary
separation}, how careful the person is), and where the token starts
(\textbf{starting point}, bias).

\paragraph{Psychology words in plain language}

\begin{itemize}
\tightlist
\item
  \textbf{Participant (subject):} a person who takes part in an
  experiment.
\item
  \textbf{Stimulus (plural stimuli):} whatever is shown or played to the
  participant, e.g.~a picture, a word or a tone.
\item
  \textbf{Task:} what the participant must do, e.g.~``press left if the
  dots move left, right if they move right''. A \textbf{two-choice task}
  has exactly two possible responses.
\item
  \textbf{Trial:} one repetition of the task: one stimulus, one
  response. An experiment usually has hundreds of trials per
  participant.
\item
  \textbf{Response time (RT, reaction time):} the time from stimulus
  onset to the key press, usually in milliseconds (ms).
\item
  \textbf{Accuracy:} the proportion of correct responses.
\item
  \textbf{Condition / manipulation:} a change the experimenter makes on
  purpose (e.g.~``be fast!'' vs.~``be accurate!'' instructions, or easy
  vs.~hard stimuli), to see how behavior changes.
\item
  \textbf{Construct:} a psychological quantity we cannot observe
  directly, e.g.~\emph{caution}, \emph{attention} or \emph{memory
  strength}. A good model parameter \emph{measures} a construct.
\item
  \textbf{Behavioral data:} what we can record: choices and RTs (not
  brain signals).
\end{itemize}

\textbf{Suggested timing (60 min):} motivation 10 · what and why 10 ·
kinds of models + Marr 15 · models as data-generating machines +
pipeline 10 · first model (random walk) 12 · notation and course
logistics 3.

\subsubsection{1. What is a cognitive model?}

A \textbf{cognitive model} is an explicit, usually mathematical or
computational, description of the processes or representations that are
assumed to produce behavior. It has three ingredients:

\begin{longtable}[]{@{}
  >{\raggedright\arraybackslash}p{(\linewidth - 4\tabcolsep) * \real{0.3333}}
  >{\raggedright\arraybackslash}p{(\linewidth - 4\tabcolsep) * \real{0.3333}}
  >{\raggedright\arraybackslash}p{(\linewidth - 4\tabcolsep) * \real{0.3333}}@{}}
\caption{\centering 1. What is a cognitive model?}\tabularnewline
\toprule\noalign{}
\begin{minipage}[b]{\linewidth}\raggedright
Ingredient
\end{minipage} & \begin{minipage}[b]{\linewidth}\raggedright
Question it answers
\end{minipage} & \begin{minipage}[b]{\linewidth}\raggedright
Random walk example
\end{minipage} \\
\midrule\noalign{}
\endfirsthead
\toprule\noalign{}
\begin{minipage}[b]{\linewidth}\raggedright
Ingredient
\end{minipage} & \begin{minipage}[b]{\linewidth}\raggedright
Question it answers
\end{minipage} & \begin{minipage}[b]{\linewidth}\raggedright
Random walk example
\end{minipage} \\
\midrule\noalign{}
\endhead
\bottomrule\noalign{}
\endlastfoot
\textbf{Assumptions / structure} & What mechanism do we posit? &
Evidence accumulates in discrete steps between two boundaries \\
\textbf{Parameters} & Which quantities are free to vary across people or
conditions? & Drift \(v\), boundary separation \(a\), start point
\(z\) \\
\textbf{Mapping to data} & How does the mechanism produce observable
responses? & The boundary reached gives the choice; the number of steps
gives the RT \\
\end{longtable}

Three related terms are easy to mix up:

\begin{itemize}
\tightlist
\item
  A \textbf{theory} is a set of verbal (often informal) claims about the
  mind.
\item
  A \textbf{model} is one particular formal instantiation of a theory.
  One theory can be formalized in many different ways.
\item
  A \textbf{statistical model} (a t-test, a regression) is also a model,
  but of the \emph{data}. Its parameters (means, slopes) usually
  describe the data rather than a cognitive process. In a cognitive
  model the parameters are meant to \textbf{measure latent psychological
  quantities}, such as memory decay, caution or attention.
\end{itemize}

George Box's often-quoted remark that ``all models are wrong, but some
are useful'' sums up the attitude of this course. We never ask whether a
model is \emph{true}. We ask whether it is \emph{useful}: does it
capture the important regularities, generalize to new data, and have
interpretable parameters?

\subsubsection{2. Why bother with formal models?}

\begin{enumerate}
\def\labelenumi{\arabic{enumi}.}
\tightlist
\item
  \textbf{Precision.} Writing the model down exposes hidden assumptions
  (as in the motivation example).
\item
  \textbf{Quantitative predictions.} Models predict \emph{how much}, not
  just \emph{whether}.
\item
  \textbf{Deriving surprising consequences.} Our intuitions about
  complex systems are poor. For example, averaging individual
  exponential learning curves produces a curve that looks like a
  \emph{power} law (Week 5).
\item
  \textbf{Measurement.} Fitted parameters are measurements of latent
  constructs. \(d'\) in signal detection theory (Week 2) measures
  sensitivity separately from response bias, which raw accuracy cannot
  do.
\item
  \textbf{Comparing theories.} Two verbal theories that ``both explain
  the data'' can be pitted against each other once they are formalized
  (Weeks 6--7, 10).
\item
  \textbf{Bridge to AI.} The same formalisms (probabilistic models,
  learning rules, neural networks) are used to build artificial agents
  (Week 15).
\end{enumerate}

\subsubsection{3. Kinds of cognitive models}

\textbf{By purpose:}

\begin{longtable}[]{@{}
  >{\raggedright\arraybackslash}p{(\linewidth - 6\tabcolsep) * \real{0.2500}}
  >{\raggedright\arraybackslash}p{(\linewidth - 6\tabcolsep) * \real{0.2500}}
  >{\raggedright\arraybackslash}p{(\linewidth - 6\tabcolsep) * \real{0.2500}}
  >{\raggedright\arraybackslash}p{(\linewidth - 6\tabcolsep) * \real{0.2500}}@{}}
\caption{\centering 3. Kinds of cognitive models}\tabularnewline
\toprule\noalign{}
\begin{minipage}[b]{\linewidth}\raggedright
Type
\end{minipage} & \begin{minipage}[b]{\linewidth}\raggedright
Question
\end{minipage} & \begin{minipage}[b]{\linewidth}\raggedright
Example
\end{minipage} & \begin{minipage}[b]{\linewidth}\raggedright
Course weeks
\end{minipage} \\
\midrule\noalign{}
\endfirsthead
\toprule\noalign{}
\begin{minipage}[b]{\linewidth}\raggedright
Type
\end{minipage} & \begin{minipage}[b]{\linewidth}\raggedright
Question
\end{minipage} & \begin{minipage}[b]{\linewidth}\raggedright
Example
\end{minipage} & \begin{minipage}[b]{\linewidth}\raggedright
Course weeks
\end{minipage} \\
\midrule\noalign{}
\endhead
\bottomrule\noalign{}
\endlastfoot
\textbf{Descriptive} & \emph{What} is the shape of the data? & Power law
of practice; exponential forgetting curve & 2--7 \\
\textbf{Process (mechanistic)} & \emph{How} does the mind produce the
behavior? & Random walk / drift-diffusion; exemplar model (GCM) & 1,
14 \\
\textbf{Normative (rational)} & What \emph{should} an ideal agent do? &
Bayesian ideal observer for cue combination & 15 \\
\end{longtable}

The boundaries are fuzzy: an SDT model is partly descriptive and partly
a process model.

\textbf{By family:}

\begin{longtable}[]{@{}
  >{\raggedright\arraybackslash}p{(\linewidth - 4\tabcolsep) * \real{0.3333}}
  >{\raggedright\arraybackslash}p{(\linewidth - 4\tabcolsep) * \real{0.3333}}
  >{\raggedright\arraybackslash}p{(\linewidth - 4\tabcolsep) * \real{0.3333}}@{}}
\caption{\centering 3. Kinds of cognitive models}\tabularnewline
\toprule\noalign{}
\begin{minipage}[b]{\linewidth}\raggedright
Family
\end{minipage} & \begin{minipage}[b]{\linewidth}\raggedright
Core idea
\end{minipage} & \begin{minipage}[b]{\linewidth}\raggedright
Example
\end{minipage} \\
\midrule\noalign{}
\endfirsthead
\toprule\noalign{}
\begin{minipage}[b]{\linewidth}\raggedright
Family
\end{minipage} & \begin{minipage}[b]{\linewidth}\raggedright
Core idea
\end{minipage} & \begin{minipage}[b]{\linewidth}\raggedright
Example
\end{minipage} \\
\midrule\noalign{}
\endhead
\bottomrule\noalign{}
\endlastfoot
Symbolic / rule-based & Cognition as manipulation of symbols and
production rules & ACT-R production systems \\
Connectionist & Distributed representations, learning by weight change &
Delta-rule networks, deep networks \\
Stochastic / dynamical & Noisy processes unfolding in time & Random
walk, drift-diffusion, Rescorla--Wagner \\
Bayesian (probabilistic) & Cognition as probabilistic inference &
Bayesian observer models \\
\end{longtable}

\begin{quote}
\textbf{Watch out:} ``Bayesian modeling'' in this course means two
different things. (a) \textbf{Bayesian models of cognition} claim that
the \emph{mind} does Bayesian inference (Week 15). (b) \textbf{Bayesian
data analysis} is how \emph{we, the scientists,} estimate the parameters
of any model, Bayesian or not (Weeks 9--13). Weeks 9--13 are mostly
about (b).
\end{quote}

\subsubsection{4. Marr's three levels of analysis}

David Marr proposed analyzing an information-processing system at three
complementary levels:

\begin{longtable}[]{@{}
  >{\raggedright\arraybackslash}p{(\linewidth - 4\tabcolsep) * \real{0.3333}}
  >{\raggedright\arraybackslash}p{(\linewidth - 4\tabcolsep) * \real{0.3333}}
  >{\raggedright\arraybackslash}p{(\linewidth - 4\tabcolsep) * \real{0.3333}}@{}}
\caption{\centering 4. Marr's three levels of analysis}\tabularnewline
\toprule\noalign{}
\begin{minipage}[b]{\linewidth}\raggedright
Level
\end{minipage} & \begin{minipage}[b]{\linewidth}\raggedright
Question
\end{minipage} & \begin{minipage}[b]{\linewidth}\raggedright
Decision example
\end{minipage} \\
\midrule\noalign{}
\endfirsthead
\toprule\noalign{}
\begin{minipage}[b]{\linewidth}\raggedright
Level
\end{minipage} & \begin{minipage}[b]{\linewidth}\raggedright
Question
\end{minipage} & \begin{minipage}[b]{\linewidth}\raggedright
Decision example
\end{minipage} \\
\midrule\noalign{}
\endhead
\bottomrule\noalign{}
\endlastfoot
\textbf{Computational} & What problem is being solved, and why? & Choose
the correct alternative while balancing speed and accuracy \\
\textbf{Algorithmic / representational} & What representations and
procedures solve it? & Accumulate noisy evidence until a boundary is
reached \\
\textbf{Implementational} & How is it realized physically? & Firing
rates of neurons ramping up to a threshold \\
\end{longtable}

Normative models usually live at the computational level, process models
at the algorithmic level, and computational neuroscience models at the
implementational level. Most models in this course are
\textbf{algorithmic-level} models.

\subsubsection{5. A model is a data-generating machine}

The single most important idea for the rest of the course: \textbf{a
model specifies a probability distribution over possible data.}

\[
\mathbf{y} \sim p(\mathbf{y} \mid \boldsymbol{\theta}, \mathcal{M})
\]

Read this as: ``under model \(\mathcal{M}\) with parameters
\(\boldsymbol{\theta}\), the data \(\mathbf{y}\) are generated from the
distribution \(p\).'' Many models split this into a deterministic core
plus noise, for example
\(y_i = f(x_i; \boldsymbol{\theta}) + \varepsilon_i\) (Week 2). Others,
like the random walk, are stochastic all the way through.

This gives two directions of reasoning:

\begin{longtable}[]{@{}
  >{\raggedright\arraybackslash}p{(\linewidth - 8\tabcolsep) * \real{0.2000}}
  >{\raggedright\arraybackslash}p{(\linewidth - 8\tabcolsep) * \real{0.2000}}
  >{\raggedright\arraybackslash}p{(\linewidth - 8\tabcolsep) * \real{0.2000}}
  >{\raggedright\arraybackslash}p{(\linewidth - 8\tabcolsep) * \real{0.2000}}
  >{\raggedright\arraybackslash}p{(\linewidth - 8\tabcolsep) * \real{0.2000}}@{}}
\caption{\centering 5. A model is a data-generating machine}\tabularnewline
\toprule\noalign{}
\begin{minipage}[b]{\linewidth}\raggedright
Direction
\end{minipage} & \begin{minipage}[b]{\linewidth}\raggedright
Given
\end{minipage} & \begin{minipage}[b]{\linewidth}\raggedright
Find
\end{minipage} & \begin{minipage}[b]{\linewidth}\raggedright
Name
\end{minipage} & \begin{minipage}[b]{\linewidth}\raggedright
Course weeks
\end{minipage} \\
\midrule\noalign{}
\endfirsthead
\toprule\noalign{}
\begin{minipage}[b]{\linewidth}\raggedright
Direction
\end{minipage} & \begin{minipage}[b]{\linewidth}\raggedright
Given
\end{minipage} & \begin{minipage}[b]{\linewidth}\raggedright
Find
\end{minipage} & \begin{minipage}[b]{\linewidth}\raggedright
Name
\end{minipage} & \begin{minipage}[b]{\linewidth}\raggedright
Course weeks
\end{minipage} \\
\midrule\noalign{}
\endhead
\bottomrule\noalign{}
\endlastfoot
\textbf{Forward} & \(\boldsymbol\theta\) & Possible data \(\mathbf{y}\)
& Simulation / prediction & 1--2 \\
\textbf{Inverse} & Observed \(\mathbf{y}\) & Plausible
\(\boldsymbol\theta\) & Parameter estimation / fitting & 3--5, 9--13 \\
\end{longtable}

The inverse problem is harder. It is what we mean by ``fitting a model
to data''. Because a model can always be run forwards, we can also test
our fitting methods: simulate data from \textbf{known} parameters
\(\boldsymbol\theta^\ast\), fit, and check whether we get
\(\boldsymbol\theta^\ast\) back (\emph{parameter recovery}, Week 7).

\paragraph{Reading the math}

\(\mathbf{y} \sim p(\mathbf{y} \mid \boldsymbol{\theta}, \mathcal{M})\)
is read: ``the data \(\mathbf{y}\) are drawn from the probability
distribution \(p\), given the parameter values \(\boldsymbol\theta\) and
the model \(\mathcal{M}\).''

\begin{itemize}
\tightlist
\item
  \(\mathbf{y}\) (bold) is a \emph{list} of observations, e.g.~the 200
  responses of one participant.
\item
  The symbol \(\sim\) means ``is randomly generated from''. It reminds
  us that the same model with the same parameters gives \emph{different}
  data each time, just as a person does not respond identically on every
  trial.
\item
  The vertical bar \(\mid\) means ``given'' or ``assuming''. Everything
  to the right of it is fixed: the model's equations (\(\mathcal{M}\))
  and the numbers we plug into them (\(\boldsymbol\theta\)).
\item
  \(\boldsymbol\theta\) (theta, bold) collects all parameters,
  e.g.~\(\boldsymbol\theta = (v, a, z)\) for the random walk.
\end{itemize}

\textbf{Example.} For a coin with \(P(\text{heads}) = 0.6\), the model
is ``each flip is heads with probability \(\theta\)'',
\(\boldsymbol\theta = 0.6\), and \(\mathbf{y}\) could be (H, T, H, H,
T). Run it again and you get another sequence. \emph{Forward} direction:
choose \(\theta\), simulate \(\mathbf{y}\). \emph{Inverse} direction
(fitting, Weeks 3--5): see \(\mathbf{y}\), find a good \(\theta\).

\subsubsection{6. The modeling pipeline (our roadmap)}

\Needspace{11\baselineskip}
% Vector rendering of the supplied week-01 Mermaid pipeline.
\begin{figure}[htbp]
\centering
\begin{tikzpicture}[>=stealth, every node/.style={draw=CMForest,rounded corners=2pt,fill=CMPale,
    text=CMInk,align=center,minimum width=3.2cm,minimum height=1.05cm,font=\small},
    every path/.style={draw=CMEmerald,thick}]
\ifCMpersian
  \node (a) at (0,0) {نظریه‌ی کلامی};
  \node (b) at (4.2,0) {مدل فرمال\\(هفته‌ی ۱\lr{--}۲)};
  \node (c) at (8.4,0) {شبیه‌سازی\\(هفته‌ی ۱\lr{--}۲)};
  \node (d) at (12.6,0) {برازش روی داده\\(هفته‌ی ۳\lr{--}۵، ۹\lr{--}۱۳)};
  \node (e) at (12.6,-2.0) {ارزیابی و مقایسه\\(هفته‌ی ۶\lr{--}۷، ۱۰)};
  \node (f) at (8.4,-2.0) {بازیابی پارامتر و مدل\\(هفته‌ی ۷)};
  \node (g) at (4.2,-2.0) {تفسیر\\(هر هفته)};
\else
  \node (a) at (0,0) {Verbal theory};
  \node (b) at (4.2,0) {Formal model\\(Weeks 1–2)};
  \node (c) at (8.4,0) {Simulate\\(Weeks 1–2)};
  \node (d) at (12.6,0) {Fit to data\\(Weeks 3–5, 9–13)};
  \node (e) at (12.6,-2.0) {Evaluate and compare\\(Weeks 6–7, 10)};
  \node (f) at (8.4,-2.0) {Recover parameters and models\\(Week 7)};
  \node (g) at (4.2,-2.0) {Interpret\\(every week)};
\fi
\path[->] (a) -- (b);
\path[->] (b) -- (c);
\path[->] (c) -- (d);
\path[->] (d) -- (e);
\path[->] (e) -- (f);
\path[->] (f) -- (g);
\path[->,dashed] (g) -- (b);
\end{tikzpicture}
\ifCMpersian\caption{چرخه‌ی مدل‌سازی شناختی، بازطراحی شده از نمودار هفته‌ی اول}\else\caption{The cognitive modeling pipeline from Week 1}\fi
\label{fig:modeling-pipeline}
\end{figure}


Every lab in this course runs some version of this loop. Today we focus
on \textbf{formalize \ensuremath{\rightarrow} simulate \ensuremath{\rightarrow} interpret}, with a small preview of
\textbf{fit}.

\subsubsection{7. Our first model: a random walk to a boundary}

\textbf{Verbal \ensuremath{\rightarrow} formal.} Evidence is a single number \(X\) that starts
at \(z\). At each time step it moves \textbf{up by 1} (evidence for the
correct response) or \textbf{down by 1} (evidence for the error
response). The walk stops when \(X\) reaches the upper boundary \(a\)
(correct response) or the lower boundary \(0\) (error).

\[
X_0 = z, \qquad X_{n} = X_{n-1} + \delta_n, \qquad
\delta_n =
\begin{cases}
+1 & \text{with probability } \tfrac{1+v}{2}\\[2pt]
-1 & \text{with probability } \tfrac{1-v}{2}
\end{cases}
\]

\begin{itemize}
\tightlist
\item
  \(v \in (-1, 1)\) is the \textbf{drift}, the average step
  \(\mathbb{E}[\delta_n] = v\). It represents the \textbf{quality of
  evidence}: easy stimuli give a large \(v\).
\item
  \(a\) is the \textbf{boundary separation}. It represents
  \textbf{caution}: a wider gap needs more evidence.
\item
  \(z\) is the \textbf{start point} (\(0 < z < a\)). With \(z = a/2\)
  the decision maker is unbiased.
\item
  The number of steps until absorption, \(T\), maps onto RT through \[
  \mathrm{RT} = t_0 + T\,\Delta t,
  \] where \(t_0\) is \textbf{non-decision time} (encoding plus motor
  execution) and \(\Delta t\) is the duration of one step.
\end{itemize}

\textbf{Analytic predictions.} Let \(r = \frac{1-v}{1+v}\). For
\(v \neq 0\), the probability of a correct response and the mean number
of steps are:

\[
P_c = \frac{1 - r^{z}}{1 - r^{a}}, \qquad
\mathbb{E}[T] = \frac{a\,P_c - z}{v}.
\]

These two equations already settle the questions from the motivation:

\begin{itemize}
\tightlist
\item
  Increasing \(a\) raises \(P_c\) \textbf{and} \(\mathbb{E}[T]\): a
  \textbf{speed--accuracy trade-off}.
\item
  Increasing \(v\) raises \(P_c\) but \textbf{lowers} \(\mathbb{E}[T]\):
  easier stimuli are both faster and more accurate.
\end{itemize}

So the two manipulations leave different signatures in the data. That is
what lets us \emph{measure} caution and evidence quality separately.

\begin{quote}
\textbf{Advanced: where the formulas come from.} \(P_c\) is the classic
\emph{gambler's ruin} result. Let \(h(x)\) be the probability of hitting
\(a\) before \(0\) when starting from \(x\). Conditioning on the first
step gives \(h(x) = \tfrac{1+v}{2}h(x+1) + \tfrac{1-v}{2}h(x-1)\) with
\(h(0)=0\) and \(h(a)=1\). Solving this difference equation gives
\(h(x) = (1-r^x)/(1-r^a)\). For \(\mathbb{E}[T]\), use \emph{Wald's
identity}: the final position satisfies
\(\mathbb{E}[X_T] = z + v\,\mathbb{E}[T]\). Since \(X_T = a\) with
probability \(P_c\) and \(X_T = 0\) otherwise,
\(\mathbb{E}[X_T] = aP_c\), and rearranging gives
\(\mathbb{E}[T] = (aP_c - z)/v\). As the step size and \(\Delta t\)
shrink together, the random walk becomes the \textbf{drift-diffusion
model} (Week 14).
\end{quote}

\begin{quote}
\textbf{Advanced: a property that turns out to be a weakness.} With an
unbiased start (\(z = a/2\)), the simple random walk predicts that
correct and error responses have \textbf{identical RT distributions}.
Empirically, errors are often \emph{slower} than correct responses when
accuracy is stressed, and sometimes \emph{faster} under speed stress.
The simple model cannot produce either pattern. Adding trial-to-trial
variability in drift or start point can (lab Exercise 5). Finding where
a model fails is exactly how models improve.
\end{quote}

\paragraph{Reading the math}

\textbf{The update rule.} \(X_n = X_{n-1} + \delta_n\) says: \emph{the
evidence now} = \emph{the evidence one step ago} + \emph{one new step}.
The step \(\delta_n\) is \(+1\) or \(-1\). The probability of \(+1\) is
\(\tfrac{1+v}{2}\). With \(v = 0\) this is \(0.5\) (a fair coin: no
information). With \(v = 0.2\) it is \(0.6\), so the walk drifts upward
on average. The average step is \((+1)(0.6) + (-1)(0.4) = 0.2 = v\),
which is why \(v\) is called the drift.

\textbf{The accuracy formula.} \(P_c = \dfrac{1 - r^{z}}{1 - r^{a}}\)
with \(r = \dfrac{1-v}{1+v}\).

\begin{itemize}
\tightlist
\item
  \(r\) is the ratio ``probability of a down step / probability of an up
  step''. With \(v = 0.2\): \(r = 0.4/0.6 = 2/3\). A small \(r\) means
  the walk almost always goes up.
\item
  \(r^{z}\) and \(r^{a}\) are small numbers when \(z\) and \(a\) are
  large, so \(P_c\) approaches 1. That is the math behind ``more caution
  (larger \(a\)) gives higher accuracy''.
\item
  You do not need to derive the formula (see the Advanced box). It is
  enough to plug in numbers and check it against a simulation, as in the
  lab.
\end{itemize}

\textbf{The time formula.} \(\mathbb{E}[T] = \dfrac{a P_c - z}{v}\).
\(\mathbb{E}[\cdot]\) means ``average over many trials''. The numerator
is how far the walk travels on average (it ends at \(a\) with
probability \(P_c\), at \(0\) otherwise, and started at \(z\)). Dividing
by the average speed \(v\) gives the average number of steps, just like
\emph{time = distance / speed}.

\textbf{RT.} \(\mathrm{RT} = t_0 + T\,\Delta t\): a fixed part \(t_0\)
(seeing the stimulus, pressing the key) plus the decision time, i.e.~the
number of steps times the duration of one step.

\paragraph{For an engineer}

You already know this model under another name.

\begin{itemize}
\tightlist
\item
  \textbf{Sequential detection.} The random walk is a discrete version
  of Wald's \textbf{sequential probability ratio test (SPRT)}:
  accumulate the log-likelihood ratio of the observations, and decide as
  soon as it crosses an upper or a lower threshold. Under some
  assumptions the SPRT gives the smallest average number of observations
  for given error probabilities, which is one reason the brain might use
  this strategy.
\item
  \textbf{Drift \ensuremath{\approx} signal-to-noise ratio.} Each step is ``signal +
  noise''. The drift \(v\) is the mean of a step and the randomness of
  the step is the noise, so \(v\) plays the role of the SNR per sample.
  A cleaner stimulus is a higher SNR.
\item
  \textbf{Boundary \ensuremath{\approx} decision threshold.} Raising the threshold lowers
  the error probability but increases the detection delay: the classic
  trade-off between error rate and latency in sequential detection. In
  psychology this is called the \textbf{speed--accuracy trade-off}.
\item
  \textbf{Start point \ensuremath{\approx} prior.} Starting closer to one boundary is like
  using a non-zero initial log-likelihood ratio, i.e.~a prior in favor
  of one hypothesis.
\item
  \textbf{Non-decision time \ensuremath{\approx} fixed processing delay.} Sensor/encoding
  latency plus actuator (motor) latency, added to every response.
\item
  \textbf{What is new for an engineer:} in engineering we \emph{design}
  the detector. Here we \emph{observe} a human detector and try to
  \textbf{infer} its parameters (SNR, threshold) from its decisions and
  decision times. That inverse problem is the subject of the whole
  course.
\end{itemize}

\subsubsection{8. Course-wide notation}

This table is used in \textbf{all} weeks. When a later week introduces a
symbol, it uses the meaning given here.

\textbf{Data and design}

\begin{longtable}[]{@{}
  >{\raggedright\arraybackslash}p{(\linewidth - 4\tabcolsep) * \real{0.3333}}
  >{\raggedright\arraybackslash}p{(\linewidth - 4\tabcolsep) * \real{0.3333}}
  >{\raggedright\arraybackslash}p{(\linewidth - 4\tabcolsep) * \real{0.3333}}@{}}
\caption{\centering 8. Course-wide notation}\tabularnewline
\toprule\noalign{}
\begin{minipage}[b]{\linewidth}\raggedright
Symbol
\end{minipage} & \begin{minipage}[b]{\linewidth}\raggedright
Meaning
\end{minipage} & \begin{minipage}[b]{\linewidth}\raggedright
First used
\end{minipage} \\
\midrule\noalign{}
\endfirsthead
\toprule\noalign{}
\begin{minipage}[b]{\linewidth}\raggedright
Symbol
\end{minipage} & \begin{minipage}[b]{\linewidth}\raggedright
Meaning
\end{minipage} & \begin{minipage}[b]{\linewidth}\raggedright
First used
\end{minipage} \\
\midrule\noalign{}
\endhead
\bottomrule\noalign{}
\endlastfoot
\(y_i\), \(\mathbf{y} = (y_1,\dots,y_N)\) & Observation \(i\); vector of
all observations & Wk 1 \\
\(N\) & Number of observations (trials) & Wk 1 \\
\(x_i\) & Experimental condition / independent variable of observation
\(i\) & Wk 2 \\
\(j = 1,\dots,J\) & Participant index; number of participants & Wk 5 \\
\end{longtable}

\textbf{Models and parameters}

\begin{longtable}[]{@{}
  >{\raggedright\arraybackslash}p{(\linewidth - 4\tabcolsep) * \real{0.3333}}
  >{\raggedright\arraybackslash}p{(\linewidth - 4\tabcolsep) * \real{0.3333}}
  >{\raggedright\arraybackslash}p{(\linewidth - 4\tabcolsep) * \real{0.3333}}@{}}
\caption{\centering 8. Course-wide notation}\tabularnewline
\toprule\noalign{}
\begin{minipage}[b]{\linewidth}\raggedright
Symbol
\end{minipage} & \begin{minipage}[b]{\linewidth}\raggedright
Meaning
\end{minipage} & \begin{minipage}[b]{\linewidth}\raggedright
First used
\end{minipage} \\
\midrule\noalign{}
\endfirsthead
\toprule\noalign{}
\begin{minipage}[b]{\linewidth}\raggedright
Symbol
\end{minipage} & \begin{minipage}[b]{\linewidth}\raggedright
Meaning
\end{minipage} & \begin{minipage}[b]{\linewidth}\raggedright
First used
\end{minipage} \\
\midrule\noalign{}
\endhead
\bottomrule\noalign{}
\endlastfoot
\(\mathcal{M}\), \(\mathcal{M}_1, \mathcal{M}_2\) & A model; competing
models & Wk 1 \\
\(\boldsymbol\theta\) & Vector of free parameters & Wk 1 \\
\(k\) & Number of free parameters & Wk 1 \\
\(\boldsymbol\theta^\ast\) & True (data-generating) parameter values in
a simulation & Wk 1 \\
\(\hat{\boldsymbol\theta}\) & Estimate of \(\boldsymbol\theta\) & Wk
3 \\
\(f(x;\boldsymbol\theta)\) & Deterministic part of a model's prediction
& Wk 2 \\
\(\hat y_i\) & Model prediction for observation \(i\) & Wk 3 \\
\(\varepsilon_i\) & Noise (error) term & Wk 2 \\
\end{longtable}

\textbf{Probability}

\begin{longtable}[]{@{}
  >{\raggedright\arraybackslash}p{(\linewidth - 2\tabcolsep) * \real{0.5000}}
  >{\raggedright\arraybackslash}p{(\linewidth - 2\tabcolsep) * \real{0.5000}}@{}}
\caption{\centering 8. Course-wide notation}\tabularnewline
\toprule\noalign{}
\begin{minipage}[b]{\linewidth}\raggedright
Symbol
\end{minipage} & \begin{minipage}[b]{\linewidth}\raggedright
Meaning
\end{minipage} \\
\midrule\noalign{}
\endfirsthead
\toprule\noalign{}
\begin{minipage}[b]{\linewidth}\raggedright
Symbol
\end{minipage} & \begin{minipage}[b]{\linewidth}\raggedright
Meaning
\end{minipage} \\
\midrule\noalign{}
\endhead
\bottomrule\noalign{}
\endlastfoot
\(P(\cdot)\) & Probability of an event \\
\(p(\cdot)\) & Probability mass or density function \\
\(y \sim \mathcal{D}\) & ``\(y\) is distributed as \(\mathcal{D}\)'' \\
\(\mathcal{N}(\mu, \sigma)\) & Normal distribution with mean \(\mu\) and
\textbf{standard deviation} \(\sigma\) (as in and ) \\
\(\mathrm{Bernoulli}(\theta)\), \(\mathrm{Binomial}(n, \theta)\) &
Bernoulli and binomial distributions with success probability
\(\theta\) \\
\(\mathrm{Beta}(a_0, b_0)\), \(\mathrm{Uniform}(l, u)\) & Beta and
uniform distributions \\
\(\Phi\), \(\Phi^{-1}\) & Standard normal CDF and its inverse \\
\(\mathbb{E}[\cdot]\), \(\mathrm{Var}[\cdot]\) & Expectation and
variance \\
\end{longtable}

\textbf{Fitting, evaluation and Bayesian inference}

\begin{longtable}[]{@{}
  >{\raggedright\arraybackslash}p{(\linewidth - 4\tabcolsep) * \real{0.3333}}
  >{\raggedright\arraybackslash}p{(\linewidth - 4\tabcolsep) * \real{0.3333}}
  >{\raggedright\arraybackslash}p{(\linewidth - 4\tabcolsep) * \real{0.3333}}@{}}
\caption{\centering 8. Course-wide notation}\tabularnewline
\toprule\noalign{}
\begin{minipage}[b]{\linewidth}\raggedright
Symbol
\end{minipage} & \begin{minipage}[b]{\linewidth}\raggedright
Meaning
\end{minipage} & \begin{minipage}[b]{\linewidth}\raggedright
First used
\end{minipage} \\
\midrule\noalign{}
\endfirsthead
\toprule\noalign{}
\begin{minipage}[b]{\linewidth}\raggedright
Symbol
\end{minipage} & \begin{minipage}[b]{\linewidth}\raggedright
Meaning
\end{minipage} & \begin{minipage}[b]{\linewidth}\raggedright
First used
\end{minipage} \\
\midrule\noalign{}
\endhead
\bottomrule\noalign{}
\endlastfoot
SSE, MSE, RMSD & Sum of squared errors, mean squared error, root mean
squared deviation & Wk 3 \\
\(L(\boldsymbol\theta) = p(\mathbf{y}\mid\boldsymbol\theta,\mathcal{M})\)
& Likelihood & Wk 4 \\
\(\ell(\boldsymbol\theta) = \ln L(\boldsymbol\theta)\); NLL \(= -\ell\)
& Log-likelihood; negative log-likelihood & Wk 4 \\
\(K\) & Number of cross-validation folds (not the same as \(k\)!) & Wk
6 \\
AIC, BIC & Akaike and Bayesian information criteria & Wk 7 \\
\(p(\boldsymbol\theta)\), \(p(\boldsymbol\theta\mid\mathbf{y})\) &
Prior; posterior & Wk 9 \\
\(p(\mathbf{y}\mid\mathcal{M})\) & Marginal likelihood (evidence) & Wk
10 \\
\(\mathrm{BF}_{12}\) & Bayes factor for \(\mathcal{M}_1\) over
\(\mathcal{M}_2\) & Wk 10 \\
\(\boldsymbol\theta^{(s)}\), \(s = 1,\dots,S\) & \(s\)-th MCMC draw;
number of draws & Wk 11 \\
\(C\) & Number of MCMC chains & Wk 12 \\
\(\hat R\), ESS & Potential scale reduction factor (R-hat); effective
sample size & Wk 12 \\
\end{longtable}

\textbf{Model-specific symbols}

\begin{longtable}[]{@{}
  >{\raggedright\arraybackslash}p{(\linewidth - 4\tabcolsep) * \real{0.3333}}
  >{\raggedright\arraybackslash}p{(\linewidth - 4\tabcolsep) * \real{0.3333}}
  >{\raggedright\arraybackslash}p{(\linewidth - 4\tabcolsep) * \real{0.3333}}@{}}
\caption{\centering 8. Course-wide notation}\tabularnewline
\toprule\noalign{}
\begin{minipage}[b]{\linewidth}\raggedright
Model
\end{minipage} & \begin{minipage}[b]{\linewidth}\raggedright
Symbols
\end{minipage} & \begin{minipage}[b]{\linewidth}\raggedright
First used
\end{minipage} \\
\midrule\noalign{}
\endfirsthead
\toprule\noalign{}
\begin{minipage}[b]{\linewidth}\raggedright
Model
\end{minipage} & \begin{minipage}[b]{\linewidth}\raggedright
Symbols
\end{minipage} & \begin{minipage}[b]{\linewidth}\raggedright
First used
\end{minipage} \\
\midrule\noalign{}
\endhead
\bottomrule\noalign{}
\endlastfoot
Random walk / drift-diffusion & \(v\) drift; \(a\) boundary separation;
\(z\) start point; \(t_0\) non-decision time; \(\Delta t\) step
duration; \(T\) number of steps & Wk 1, 14 \\
Signal detection theory & \(H\) hit rate; \(F\) false-alarm rate; \(d'\)
sensitivity; \(c\) criterion & Wk 2 \\
Forgetting curves & \(t\) retention interval; \(A\) initial level;
\(\lambda\) exponential decay rate; \(d\) power-law decay exponent & Wk
2 \\
GCM (categorization) & \(w_m\) attention weights; \(c\) sensitivity & Wk
14 \\
Reinforcement learning & \(\alpha\) learning rate; \(\beta\) inverse
temperature; \(Q_t(o)\) value of option \(o\) on trial \(t\) & Wk 15 \\
\end{longtable}

\textbf{Overloaded symbols (kept because the literature uses them):}
\(c\) is the SDT criterion except in Week 14, where it is the GCM
sensitivity. \(t\) is the retention interval in forgetting models and
the trial index in RL. \(\alpha\) is only used as the RL learning rate.
We never use \(\alpha\) for a significance level.

\paragraph{For the instructor}

\begin{itemize}
\tightlist
\item
  \textbf{Start with the questions, not the model.} Read the verbal
  theory aloud and ask students to answer the four questions of the
  Motivation section. They will disagree. That disagreement is the
  reason formal models exist.
\item
  \textbf{Act out the random walk.} Draw a vertical line on the board
  with a top line (``correct'') and a bottom line (``error''). Let a
  student roll a die: 1--4 means up, 5--6 means down
  (\(v \approx 0.33\)). Count steps until a line is reached. Do this
  three times and you have three trials with different RTs. Then move
  the lines further apart and repeat.
\item
  \textbf{Keep the math light in this week.} Students only need to
  \emph{use} the formulas for \(P_c\) and \(\mathbb{E}[T]\), not derive
  them. Show that the numbers match the simulation in the lab: that is
  the main message (simulation and analysis agree).
\item
  \textbf{Typical student questions.}

  \begin{itemize}
  \tightlist
  \item
    \emph{``Is evidence really discrete steps?''} No.~It is a
    simplification; Week 14 shows the continuous version
    (drift-diffusion), which gives similar predictions.
  \item
    \emph{``Why not just compare mean RTs with a t-test?''} A t-test
    says \emph{that} groups differ; the model says \emph{how} (caution
    vs.~evidence quality).
  \item
    \emph{``Which model is the true one?''} None. All models are
    simplifications; we ask which one is useful (Box's quote).
  \end{itemize}
\item
  \textbf{Time management.} Sections 3--4 (model types, Marr) can be
  shortened to 10 minutes if needed. Do not skip Section 7: everything
  in the lab depends on it.
\end{itemize}

\begin{center}\rule{0.5\linewidth}{0.5pt}\end{center}

\subsection{Worked example}

\textbf{Setup.} An unbiased decision maker with drift \(v = 0.2\),
boundary separation \(a = 6\) and start point \(z = 3\). Non-decision
time \(t_0 = 300\) ms; each step lasts \(\Delta t = 20\) ms.

\textbf{Step 1: step probabilities.}
\(P(\delta = +1) = (1 + 0.2)/2 = 0.6\) and \(P(\delta = -1) = 0.4\).

\textbf{Step 2: simulate one trial by hand.} Draw uniform random numbers
\(u_n\). Step up if \(u_n < 0.6\), otherwise step down.

\begin{longtable}[]{@{}llll@{}}
\caption{\centering Worked example}\tabularnewline
\toprule\noalign{}
Step \(n\) & \(u_n\) & \(\delta_n\) & \(X_n\) \\
\midrule\noalign{}
\endfirsthead
\toprule\noalign{}
Step \(n\) & \(u_n\) & \(\delta_n\) & \(X_n\) \\
\midrule\noalign{}
\endhead
\bottomrule\noalign{}
\endlastfoot
0 & --- & --- & 3 \\
1 & 0.31 & +1 & 4 \\
2 & 0.72 & −1 & 3 \\
3 & 0.55 & +1 & 4 \\
4 & 0.12 & +1 & 5 \\
5 & 0.48 & +1 & \textbf{6 = a} \ensuremath{\rightarrow} stop \\
\end{longtable}

The walk hits the upper boundary: a \textbf{correct} response after
\(T = 5\) steps, so \(\mathrm{RT} = 300 + 5 \times 20 = 400\) ms.

\textbf{Step 3: analytic accuracy.} \(r = (1-0.2)/(1+0.2) = 2/3\), so
\(r^3 = 8/27 \approx 0.296\) and \(r^6 = 64/729 \approx 0.088\).

\[
P_c = \frac{1 - 8/27}{1 - 64/729} = \frac{1}{1 + 8/27} = \frac{27}{35} \approx 0.771.
\]

(When \(z = a/2\), the formula simplifies to \(P_c = 1/(1 + r^{a/2})\).)

\textbf{Step 4: analytic mean RT.}

\[
\mathbb{E}[T] = \frac{6 \times 0.771 - 3}{0.2} = \frac{57/35}{0.2} \approx 8.14 \text{ steps},
\qquad
\mathbb{E}[\mathrm{RT}] = 300 + 8.14 \times 20 \approx 462.9 \text{ ms}.
\]

\textbf{Step 5: two manipulations, two signatures.}

\begin{longtable}[]{@{}
  >{\raggedright\arraybackslash}p{(\linewidth - 10\tabcolsep) * \real{0.1667}}
  >{\raggedright\arraybackslash}p{(\linewidth - 10\tabcolsep) * \real{0.1667}}
  >{\raggedright\arraybackslash}p{(\linewidth - 10\tabcolsep) * \real{0.1667}}
  >{\raggedright\arraybackslash}p{(\linewidth - 10\tabcolsep) * \real{0.1667}}
  >{\raggedright\arraybackslash}p{(\linewidth - 10\tabcolsep) * \real{0.1667}}
  >{\raggedright\arraybackslash}p{(\linewidth - 10\tabcolsep) * \real{0.1667}}@{}}
\caption{\centering Worked example}\tabularnewline
\toprule\noalign{}
\begin{minipage}[b]{\linewidth}\raggedright
Condition
\end{minipage} & \begin{minipage}[b]{\linewidth}\raggedright
\(v\)
\end{minipage} & \begin{minipage}[b]{\linewidth}\raggedright
\(a\) (\(z = a/2\))
\end{minipage} & \begin{minipage}[b]{\linewidth}\raggedright
\(P_c\)
\end{minipage} & \begin{minipage}[b]{\linewidth}\raggedright
\(\mathbb{E}[T]\)
\end{minipage} & \begin{minipage}[b]{\linewidth}\raggedright
\(\mathbb{E}[\mathrm{RT}]\) (ms)
\end{minipage} \\
\midrule\noalign{}
\endfirsthead
\toprule\noalign{}
\begin{minipage}[b]{\linewidth}\raggedright
Condition
\end{minipage} & \begin{minipage}[b]{\linewidth}\raggedright
\(v\)
\end{minipage} & \begin{minipage}[b]{\linewidth}\raggedright
\(a\) (\(z = a/2\))
\end{minipage} & \begin{minipage}[b]{\linewidth}\raggedright
\(P_c\)
\end{minipage} & \begin{minipage}[b]{\linewidth}\raggedright
\(\mathbb{E}[T]\)
\end{minipage} & \begin{minipage}[b]{\linewidth}\raggedright
\(\mathbb{E}[\mathrm{RT}]\) (ms)
\end{minipage} \\
\midrule\noalign{}
\endhead
\bottomrule\noalign{}
\endlastfoot
Baseline & 0.2 & 6 & 0.771 & 8.14 & 462.9 \\
``Be accurate!'' (more caution) & 0.2 & 10 & 0.884 & 19.18 & 683.6 \\
Easier stimulus (better evidence) & 0.3 & 6 & 0.865 & 7.30 & 446.0 \\
\end{longtable}

\textbf{Interpretation.} Accuracy instructions buy about 11 percentage
points of accuracy at a cost of roughly 220 ms. Easier stimuli give a
\emph{similar} accuracy gain while making responses slightly
\emph{faster}. A researcher who looked only at accuracy could not tell
the two situations apart. The model, by using accuracy \emph{and} RT
together, attributes them to different psychological causes (caution
vs.~evidence quality).

\begin{center}\rule{0.5\linewidth}{0.5pt}\end{center}

\Needspace{11\baselineskip}

\subsection{Lab (2 hours)}

\subsubsection{Goals}

\begin{enumerate}
\def\labelenumi{\arabic{enumi}.}
\tightlist
\item
  Formulate and evaluate the random walk model, simulate an experiment, and
  \textbf{verify} the simulation against the analytic formulas.
\item
  Use the model to \textbf{interpret} a speed--accuracy trade-off.
\item
  Get a first taste of \textbf{fitting}: recover the hidden parameters
  of a ``mystery participant'', \textbf{evaluate} the fit, and
  \textbf{interpret} the result.
\end{enumerate}

\textbf{Timing:} Part A 35 min · Part B 15 min · Part C 30 min · Part D
15 min · Part E 25 min. Exercises not finished in the lab are finished
at home.

\subsubsection{Part C: Formulate, simulate and verify the random walk
(30 min)}

First, look at what the model does: plot a few evidence trajectories.

Now simulate a whole experiment and compare the simulation with the
analytic predictions. \textbf{Always check a simulation against
something you know}; this habit will save you many times this term.

\textbf{What to notice:} - The simulation matches the analytic values to
within sampling noise. - The RT distributions are \textbf{right-skewed}
(a long tail of slow responses), which is the typical shape of empirical
RT distributions. - Correct and error RTs are almost identical (see the
Advanced box in Lecture §7).

\subsubsection{Part D: Interpret, the speed--accuracy trade-off (15
min)}

Moving along one curve (increasing \(a\)) trades speed for accuracy.
Moving \emph{between} curves (increasing \(v\)) improves accuracy at the
same or lower RT. The table reproduces the worked example: e.g.,
\(v = 0.2, a = 10\) gives accuracy 0.884 and mean RT 683.6 ms.

\textbf{Discuss with your neighbor:} an experimenter finds that a drug
makes participants slower \emph{and} more accurate. Which parameter
would you suspect changed? What if the drug made them slower and
\emph{less} accurate?

\subsubsection{Part E: A first fit, evaluation and interpretation (25
min)}

\textbf{Don't peek}: try to recover them.

We assume \(t_0 = 300\) ms and \(\Delta t = 20\) ms are known and
\(z = a/2\). Only \(v\) and \(a\) are free (\(k = 2\)).

\textbf{Fitting by grid search.} For every combination of \(v\) and
\(a\) on a grid, compute the predicted accuracy and mean RT and measure
how far they are from the observed values. Each difference is expressed
in standard-error units so that accuracy and RT are on a comparable
scale. Week 3 develops this idea properly.

\textbf{Evaluating the fit.} We fitted only two summary statistics
(accuracy and mean RT). A stronger test is whether the model also
reproduces something we did \emph{not} fit: the \textbf{shape} of the RT
distribution. We compare RT quantiles of correct responses in the data
with those of a large simulation at the best-fitting parameters.

The hidden values were \(v^\ast = 0.16\) and \(a^\ast = 10\). The grid
search recovered \(a\) exactly and \(v\) to within one grid step (0.15).
The small miss is not a bug: with 1000 trials, the observed accuracy and
mean RT differ from their true values by sampling noise, so the
best-fitting parameters do too. (Week 5 shows how to put an uncertainty
interval around estimates like these.) RTs are multiples of 20 ms, which
is why many quantiles match exactly. The model reproduces the shape of
the distribution well, with only a slight underestimate of the slow tail
(0.9 quantile).

In psychological terms, what do the recovered \(v\) and \(a\) say about
this participant? 2. The fit used only two numbers from the data. Why is
the quantile comparison a more convincing check than the discrepancy
value itself? 3. If this participant were re-tested under ``be as fast
as possible'' instructions, which parameter should change and which
should stay the same, according to the model?

\subsubsection{Exercises (increasing difficulty)}

\textbf{Exercise 1 \ensuremath{\star} (sampling noise).} Simulate the baseline model
(\(v = 0.2\), \(a = 6\)) with \(N = 20\), \(100\) and \(1000\) trials,
five different seeds each. Make a table of the observed accuracies.
Compare their spread with the standard error \(\sqrt{P_c(1-P_c)/N}\).
\emph{Expected:} with \(N = 20\) accuracies can easily range from about
0.6 to 0.95. The spread shrinks roughly with \(1/\sqrt{N}\).

\textbf{Exercise 2 \ensuremath{\star} (two conditions).} Simulate a ``speed'' condition
(\(a = 4\)) and an ``accuracy'' condition (\(a = 10\)) with the same
\(v = 0.2\), 1000 trials each. Report accuracy and mean RT per condition
and plot both RT distributions. Which aspect of the RT distribution
changes most: its minimum, its mode, or its tail?

\textbf{Exercise 3 \ensuremath{\star}\ensuremath{\star} (response bias).} Set \(a = 8\) and \(z = 5\)
(start closer to the correct boundary). Compute \(P_c\) analytically and
by simulation. Then compare the mean RT of correct and error responses.
Explain the result in terms of the distance each boundary is from the
start point. What would \(z = 3\) do?

\textbf{Exercise 4 \ensuremath{\star}\ensuremath{\star} (checking a model property).} For \(z = a/2\),
verify by simulation (20 000 trials) that correct and error RT
\emph{quantiles} are equal to within sampling noise. Empirical data
often show slower errors in difficult, accuracy-stressed conditions. Is
the simple random walk falsified by such data? What would you change?

\(v_i \sim \mathcal{N}(v, \eta)\) Use \(v = 0.2\), \(\eta = 0.2\),
\(a = 8\). Show that \textbf{errors are now slower than correct
responses} (expected with 20 000 trials: accuracy \ensuremath{\approx} 0.75, about 12 steps
for correct vs.~14 steps for error responses), and explain why.
\emph{Hint:} which trials are most likely to end in an error, and how
long do those trials take?

\begin{center}\rule{0.5\linewidth}{0.5pt}\end{center}

\subsection{Connection to AI \& LLMs}

\textbf{Marr's levels for a chatbot.} A large language model (LLM) can
be described at all three of Marr's levels, and the levels can come
apart just as they do for people. \emph{Computational:} the model is
trained to predict the next word (token) in a text. \emph{Algorithmic:}
a transformer network computes a score (a ``logit'') for every possible
next token and turns the scores into probabilities.
\emph{Implementational:} matrix multiplications on GPUs. A claim like
``the LLM understands the question'' is usually made at the
computational level, while the evidence (its outputs) is produced by the
algorithm. Keep the levels apart when you read such claims.

\textbf{A familiar decision rule.} The final step of an LLM is a choice
rule with a noise parameter, just like the cognitive models in this
course. The probability of token \(i\) is a softmax with
\textbf{temperature} \(T\): \(P(i) = e^{z_i/T}/\sum_j e^{z_j/T}\). A low
\(T\) gives near-deterministic choices and a high \(T\) gives more
random ones. This is the same rule as the softmax in Week 15, with
\(\beta = 1/T\).

- \ensuremath{\star} Is the LLM's temperature closer to the random walk's \emph{noise} or
to its \emph{boundary} (Week 1, Part C)? Explain. - \ensuremath{\star}\ensuremath{\star} Discussion:
someone says ``the chatbot gave a different answer when I asked again,
so it does not know the answer.'' Use this box to explain what else
could be going on.

\begin{center}\rule{0.5\linewidth}{0.5pt}\end{center}

\subsection{Common mistakes \& misconceptions}

\begin{enumerate}
\def\labelenumi{\arabic{enumi}.}
\tightlist
\item
  \textbf{``A model that fits well is true.''} Good fit is necessary but
  not sufficient. A flexible model can fit almost anything, including
  noise (Week 6). Fit says little until you know what the model
  \emph{could not} have fitted.
\item
  \textbf{``More parameters make a better model.''} More parameters make
  a model \emph{more flexible}, which is not the same as better. We will
  penalize flexibility explicitly (AIC/BIC, cross-validation, Bayes
  factors).
\item
  \textbf{Confusing simulation with fitting.} Simulation goes from
  parameters to data; fitting goes from data to parameters. Students
  often ``fit'' by tweaking parameters until a plot looks nice. That is
  a legitimate \emph{exploration} but not an estimation method.
\item
  \textbf{Treating parameters as directly observable.} \(v\) and \(a\)
  are \emph{inferred} quantities whose values depend on the model's
  assumptions (e.g., that \(t_0\) is known). Change the model and the
  ``same'' parameter may change meaning.
\item
  \textbf{Not verifying simulations.} A bug in a simulator produces
  wrong conclusions that \emph{look} plausible. Compare against an
  analytic result, a limiting case (e.g., \(v \to 0\) should give
  \(P_c \to z/a\)), or a published value.
\item
  \textbf{Mixing up the two meanings of ``Bayesian''.} Bayesian models
  of cognition (a claim about the mind) vs.~Bayesian data analysis (a
  method for scientists).
\item
  \textbf{Treating Marr's levels as competing.} They are complementary
  descriptions of the same system, not rival theories.
\end{enumerate}

\begin{center}\rule{0.5\linewidth}{0.5pt}\end{center}

\subsection{Check your understanding}

\textbf{1.} In one sentence each, define a \emph{theory}, a
\emph{cognitive model} and a \emph{statistical model}. \textbf{Answer} A
theory is a set of (usually verbal) claims about mental processes. A
cognitive model is a formal, explicit instantiation of such claims that
can generate predictions, with parameters that represent psychological
quantities. A statistical model describes the structure of the data
(e.g., group means, slopes) without committing to a cognitive mechanism.

\textbf{2.} Classify each as descriptive, process or normative: (a)
\(\mathrm{RT} = A\,N^{-b}\) as a function of practice trials \(N\); (b)
the random walk model; (c) an ideal observer that combines visual and
haptic cues weighted by their reliabilities. \textbf{Answer} (a)
Descriptive: it summarizes the shape of the data without saying why. (b)
Process: it specifies a mechanism that produces choices and RTs. (c)
Normative: it specifies optimal behavior given the available
information.

\textbf{3.} At which of Marr's levels does the random walk model mainly
live? Give an example of a claim about the same decision at the
implementational level. \textbf{Answer} The algorithmic/representational
level: it describes a procedure (accumulate evidence to a boundary). An
implementational claim would be, e.g., ``the accumulated evidence is
carried by the firing rate of neurons in a specific brain area, and the
response is triggered when that rate reaches a threshold.''

\textbf{4.} What is the difference between the forward and the inverse
problem? Which one did we solve in Part E of the lab? \textbf{Answer}
Forward: given parameters, generate (predict) data, i.e., simulation.
Inverse: given data, infer the parameters, i.e., fitting. Part E solved
an inverse problem (by grid search), but it relied on solving the
forward problem many times, once per grid point.

\textbf{5.} Using the formulas, what happens to \(P_c\) when \(v = 0\)?
(Hint: think about the limit, or about the symmetry of the walk.)
\textbf{Answer} With no drift the walk is symmetric, and the probability
of hitting the upper boundary equals \(z/a\) (the limit of the formula
as \(v \to 0\)). With an unbiased start, \(z = a/2\), accuracy is 0.5,
i.e., chance. That is sensible: with no evidence the decision maker is
guessing.

\textbf{6.} Two groups differ in accuracy (0.80 vs.~0.90). Why can't we
tell from accuracy alone whether they differ in caution or in evidence
quality? What extra data resolve this? \textbf{Answer} Both a larger
boundary \(a\) and a larger drift \(v\) raise accuracy, so accuracy
alone confounds them. RT resolves it: higher caution makes responses
slower, while better evidence makes them faster (worked example, Step
5). Fitting the model to accuracy and RT jointly separates the two.

\textbf{7.} Why did we simulate data from \emph{known} parameters before
trusting the fitting procedure? \textbf{Answer} Only with known true
values can we check whether the method returns the right answer
(parameter recovery). If a method cannot recover parameters from data
generated by the model itself, its estimates from real data are not
interpretable. This is developed fully in Week 7.

\begin{center}\rule{0.5\linewidth}{0.5pt}\end{center}

\Needspace{11\baselineskip}

\subsection{Assignment}

\subsubsection{A0: Verbal theory audit}

A0 does not carry points, but it is required: an incomplete A0 must be
revised before the Week 2 lab check-off can be awarded.

\textbf{Part 2: Verbal theory audit (max. 1 page).} Choose one verbal
theory from your earlier courses, for example ``levels of processing'',
``attention as a spotlight'', ``practice makes perfect'', ``dual-process
theories of reasoning'' or ``interference causes forgetting''.

\begin{enumerate}
\def\labelenumi{\arabic{enumi}.}
\tightlist
\item
  State the theory in 1--2 sentences.
\item
  List \textbf{at least five} decisions the theory leaves unspecified
  that a formal model would have to make (e.g., functional form, what is
  noisy, what the free parameters are, how internal states map onto the
  measured behavior).
\item
  For \textbf{two} of these decisions, propose a concrete formal choice
  (an equation or an algorithmic step).
\item
  Describe \textbf{one} prediction that would differ depending on which
  choice is made, and an experiment that could tell them apart.
\end{enumerate}

\textbf{Feedback criteria (pass / revise):}

\begin{longtable}[]{@{}
  >{\raggedright\arraybackslash}p{(\linewidth - 4\tabcolsep) * \real{0.3333}}
  >{\raggedright\arraybackslash}p{(\linewidth - 4\tabcolsep) * \real{0.3333}}
  >{\raggedright\arraybackslash}p{(\linewidth - 4\tabcolsep) * \real{0.3333}}@{}}
\caption{\centering A0: Verbal theory audit}\tabularnewline
\toprule\noalign{}
\begin{minipage}[b]{\linewidth}\raggedright
Criterion
\end{minipage} & \begin{minipage}[b]{\linewidth}\raggedright
Pass
\end{minipage} & \begin{minipage}[b]{\linewidth}\raggedright
Revise
\end{minipage} \\
\midrule\noalign{}
\endfirsthead
\toprule\noalign{}
\begin{minipage}[b]{\linewidth}\raggedright
Criterion
\end{minipage} & \begin{minipage}[b]{\linewidth}\raggedright
Pass
\end{minipage} & \begin{minipage}[b]{\linewidth}\raggedright
Revise
\end{minipage} \\
\midrule\noalign{}
\endhead
\bottomrule\noalign{}
\endlastfoot
Underspecification & \ensuremath{\geq} 5 distinct, genuinely unspecified decisions &
Fewer than 5, or restatements of the theory \\
Formalization & Two concrete, unambiguous formal choices & Vague (``use
a function'') \\
Diverging prediction & A specific prediction plus an experiment that
distinguishes the choices & No prediction, or a prediction that does not
depend on the choice \\
\end{longtable}

The best audits will be discussed briefly at the start of the Week 2
lecture.

\begin{center}\rule{0.5\linewidth}{0.5pt}\end{center}

\subsection{Further reading}

\begin{itemize}
\tightlist
\item
  \textbf{Farrell \& Lewandowsky (2018), chapters 1--2.} Introduction to
  computational modeling and the move from verbal theories to formal
  models. \emph{(Main reading for this week.)}
\item
  \textbf{Verguts (2022), introductory chapter {[}verify{]}.} A short,
  accessible motivation for modeling cognitive processes.
\item
  \textbf{Anderson (2014), introductory chapter {[}verify{]}.} Places
  cognitive modeling alongside computational neuroscience, useful
  context for Marr's implementational level.
\item
  \emph{Optional classic (outside the course list):} Marr, D. (1982).
  \emph{Vision.} W. H. Freeman. The original source of the three levels
  of analysis.
\end{itemize}

Assignment A1 will be released.
